На основе представления молекул в формате SMILES было разработано программное обеспечение для распознавания химических структур по изображениям. Для решения данной задачи применяются методы OCSR. 
В качестве основы для извлечения SMILES из изображений химических формул были рассмотрены три модели, представляющие различные подходы к задаче OCSR:

% Следует отметить, что OCSR представляет собой не отдельный тип модели, а прикладную задачу, для решения которой используются методы компьютерного зрения и глубокого обучения. В контексте классификации искусственного интеллекта такие методы относятся к области глубокого обучения, представленной на рис.~\ref{fig:classifications_ai} как наиболее внутренний уровень иерархии.

\begin{enumerate}
	\item \textbf{MolScribe.}  
	MolScribe представляет собой структурно-ориентированный подход к распознаванию химических изображений. Ключевая идея модели заключается в том, что химическую структуру целесообразно рассматривать не как набор независимых символов или линий, а как молекулярный граф, который затем может быть преобразован в строковое представление. В рамках данного подхода изображение сначала переводится в набор визуальных признаков, после чего на их основе восстанавливается структура молекулы~\cite{refMolScribe}. Таким образом, модель определяет положение атомов, типы связей между ними, наличие ароматичности, зарядов и особенности стереохимии. Благодаря этому MolScribe можно отнести к моделям, ориентированным на восстановление химического графа, а не только на последовательное распознавание элементов изображения.
	
	\item \textbf{DECIMER (Deep Learning for Chemical Image Recognition).}  
	DECIMER представляет собой модель OCSR, основанную на методах глубокого обучения и обученную на больших объёмах синтетически сгенерированных данных~\cite{refDECIMER}. Архитектурно данный подход ориентирован на преобразование изображения химической структуры в строковое представление молекулы. В основе модели лежит схема, в которой сначала из изображения извлекаются визуальные признаки, а затем они используются для генерации последовательности токенов, описывающих химическую структуру. В качестве текстового представления в DECIMER применяется SELFIES, что позволяет повысить устойчивость генерации по сравнению с непосредственным предсказанием SMILES~\cite{refDECIMERJournal}. Существенным достоинством DECIMER является высокая точность распознавания на больших наборах данных, однако по сравнению с графо-ориентированными подходами модель в меньшей степени отражает молекулу именно как графовую структуру, поскольку её основной выход формируется в виде последовательности символов~\cite{refDECIMERJournal}.
	
	\item \textbf{Qwen3-VL.}  
	Qwen3-VL представляет собой мультимодальную vision-language модель, способную анализировать графическую информацию и формировать текстовое описание содержимого изображения. В контексте задачи OCSR такая модель может рассматриваться как универсальный инструмент визуально-текстового преобразования, потенциально пригодный для интерпретации изображений химических структур. В отличие от специализированных моделей, таких как MolScribe и DECIMER, Qwen3-VL не разрабатывалась исключительно для задач химического распознавания, однако её мультимодальная архитектура позволяет использовать её в качестве сравнительного или вспомогательного подхода при анализе химических изображений~\cite{refQwen3VL}.
\end{enumerate}

Для применения локальных моделей MolScribe и DECIMER в работе использовались их встроенные механизмы обработки изображений химических структур. В случае Qwen3-VL он был также развёрнут локально, но взаимодействие с моделью осуществлялось программно через API с использованием системного промпта и набора демонстрационных примеров. Такой подход соответствует парадигме few-shot in-context learning, при которой модель адаптируется к требуемому формату ответа на основе контекста запроса без дополнительного дообучения параметров модели.

Концептуально использование метода <<few-shot in-context learning>> в молекулярных задачах согласуется с результатами работы C.~Fifty, J.~Leskovec и S.~Thrun~\cite{refCAMP}, где предложен алгоритм CAMP для few-shot предсказания молекулярных свойств. В данной статье показано, что модель может выполнять предсказание на основе контекста из пар <<объект--метка>> за один проход, без этапа тонкой настройки модели, и при малом числе примеров превосходит ряд классических мета-обучающих-подходов~\cite{refCAMP}. Хотя работа~\cite{refCAMP} посвящена не распознаванию химических структур по изображениям, а задаче предсказания молекулярных свойств, она подтверждает применимость самой идеи few-shot in-context learning в химических задачах при ограниченном числе демонстраций. Поэтому в настоящей работе данный принцип используется как методологическое основание для применения Qwen3-VL в задаче распознавания химических формул по изображениям: модели предоставляется системная инструкция и несколько примеров соответствия между изображением структуры и требуемым текстовым представлением.
